% Figure 11.2: Wick rotation contour.
% Source: physical PDF p. 499 (printed p. 476).
\begingroup
\renewcommand{\thefigure}{11.\arabic{figure}}
\setcounter{figure}{1}
\begin{figure}[tbp]
\centering
\begin{tikzpicture}[x=1cm,y=1cm,every node/.style={font=\small}]
  \draw[line width=0.65pt] (-3.45,0) -- (3.45,0);
  \draw[line width=0.65pt] (0,-3.25) -- (0,3.45);
  \node at (-2.75,0.42) {$\mathsf{x}$};
  \node at (2.75,-0.42) {$\mathsf{x}$};
  \draw[-{Stealth[length=5pt,width=3.8pt]},line width=0.75pt]
    (1.92,0.45) .. controls (1.65,1.35) and (1.12,2.03) .. (0.70,2.13);
\end{tikzpicture}
\caption{Wick rotation of the $p^0$ contour of integration. Small x's mark the poles in the $p^0$ complex plane; the arrow indicates the direction of rotation of the contour of integration, from the real to the imaginary $p^0$-axes.}
\label{fig:11.2}
\end{figure}
\endgroup
